% META: {"id":"1SPE-GEOREP-EX-015","chapitre":"1SPE-GEOMETRIE-REPEREE","type_objet":"exercice","capacites":["1SPE-GEOMETRIE-REPEREE-C2"],"capacites_codes":["C2"],"methodes":["M2"],"parcours":2,"competences":["calculer","raisonner"],"duree_min":12,"mode_creation":"ex_nihilo","sources_inspiration":[],"fichier_tex":"chapitres/1SPE-GEOMETRIE-REPEREE/exercices/1SPE-GEOREP-EX-015.tex","corrige_tex":"chapitres/1SPE-GEOMETRIE-REPEREE/corriges/1SPE-GEOREP-CO-015.tex","status":"approved","origin":"generated","status_history":["generated","audited","accepted","approved"],"release_acceptance":"RELEASE_OWNER_BATCH_ACCEPTANCE","acceptance_closure_digest":"sha256:9b3ccf9a81c5520fbb7e03b7b2d3e7bf2057a02834b6d908bf3be5f49f3b553f","acceptance_receipt_digest":"sha256:4fad86f0282e0fa4e0419cca1ad6379ae7af5a280c04a4a90880f90a1c044242"}
% BEGIN-VERIFY
% from sympy import *
% x, y = symbols('x y')
% # Droite perpendiculaire a D: 2x - 5y + 3 = 0 passant par P(4,1)
% # Vecteur directeur de D: (5,2), qui est normal de la perpendiculaire
% # 5(x-4) + 2(y-1) = 0 => 5x + 2y - 22 = 0
% assert 5*4 + 2*1 - 22 == 0
% assert 2*5 + (-5)*2 == 0
% END-VERIFY
\begin{exercice}{1SPE-GEOREP-EX-015}{2}{12}
Déterminer l'équation de la droite $\mathcal{D}'$ perpendiculaire à $\mathcal{D} : 2x - 5y + 3 = 0$ passant par $P(4 ; 1)$.
\end{exercice}
